\documentclass[conference]{IEEEtran}
\IEEEoverridecommandlockouts
% The preceding line is only needed to identify funding in the first footnote. If that is unneeded, please comment it out.
\usepackage{cite}
\usepackage{amsmath,amssymb,amsfonts}
\usepackage{algorithmic}
\usepackage{graphicx}
\usepackage{textcomp}
\usepackage{xcolor}
\usepackage{multirow}
\usepackage[T5]{fontenc}
\usepackage[utf8]{inputenc}
\usepackage{tabularx}
\usepackage{array}
\usepackage{booktabs}
\usepackage{algorithm}
\usepackage{hyperref}

\def\BibTeX{{\rm B\kern-.05em{\sc i\kern-.025em b}\kern-.08em
    T\kern-.1667em\lower.7ex\hbox{E}\kern-.125emX}}
\begin{document}

% -----------------------------------Title-----------------------------------
\title{
    A Comparative Study of CNN-LSTM Architectures for Visual Question Answering\\
    % {\footnotesize \textsuperscript{*}Note: Sub-titles are not captured in Xplore and
    % should not be used}
    % \thanks{Identify applicable funding agency here. If none, delete this.}
}

% -----------------------------------Author-----------------------------------
\author{
    \IEEEauthorblockN{
        Quang Huy Nguyen\textsuperscript{*},
        Tri Hieu Phan\textsuperscript{*},
        Huu Thinh Kieu\textsuperscript{*}
    }
    \IEEEauthorblockA{
        \textit{Faculty of Information Technology, Ton Duc Thang University} \\
        Ho Chi Minh City, Vietnam \\
        \{523h0140, 523h0133, 512h0180\}@student.tdtu.edu.vn
    }
    \\
    \IEEEauthorblockN{
        Anh Cuong Le \textsuperscript{**}
    }
    \IEEEauthorblockA{
        \textit{NLP \& KD Laboratory, Faculty of Information Technology, Ton Duc Thang University} \\
        Ho Chi Minh City, Vietnam \\
        leanhcuong@tdtu.edu.vn
    }
    \thanks{*These authors contributed equally to this work. **This is our lecturer.}
}

\maketitle

% -----------------------------------Abstract-----------------------------------
\begin{abstract}

\end{abstract}

% -----------------------------------Keyword-----------------------------------
\begin{IEEEkeywords}

\end{IEEEkeywords}

% -----------------------------------Introduction-----------------------------------
\section{\textbf{Introduction}}

% -----------------------------------Related Work-----------------------------------
\section{\textbf{Related Work}}

Visual Question Answering (VQA) sits at the intersection of computer vision and natural language processing, requiring a system to jointly understand an image and a textual question to produce a correct answer. Research in this space has driven advances in visual encoding, sequence modeling, cross-modal fusion, and answer generation. This section reviews the key prior works underlying the components of our proposed system.

Convolutional Neural Networks (CNNs) are the standard backbone for visual feature extraction in VQA systems. He et al.~\cite{Kaiming_deep_res} proposed residual learning, which reformulates each layer as learning a residual function with respect to its input rather than an unreferenced mapping. This design allows networks to be trained to much greater depth without degradation, and ResNet variants consistently achieve strong performance across visual recognition benchmarks. We adopt pretrained ResNet-18 and ResNet-50 as our image encoders, retaining spatial feature maps from the final convolutional layer rather than the global pooled vector, so that positional information across image regions is preserved for the subsequent attention mechanism.

% -----------------------------------Model Architecture-----------------------------------
\section{\textbf{Model Architecture}}

\begin{figure*}[t]
    \centering
    \includegraphics[width=1.0\linewidth]{figure/cnn_feature_extraction.png}
    \caption{CNN feature extraction pipeline. An input infographic is encoded by ResNet-50 into a spatial feature tensor and then reshaped into regional vectors for downstream attention.}
    \label{fig:cnn_feature_extraction}
\end{figure*}


% -----------------------------------Experiments-----------------------------------
\section{\textbf{Experiments}}


\begin{table}[t]
    \centering
    % \small
    \caption{Dataset statistics and splits of ViInfographicVQA.}
    \label{tab:dataset_stats}
    % \footnotesize

    \begin{tabular}{lccc}
        \toprule
        \textbf{Subset} & \textbf{Images} & \textbf{Groups} & \textbf{QAs} \\
        \midrule
        Single-Image (train) & 1788 & -- & 12521 \\
        Single-Image (test)  & 192  & -- & 1374  \\
        Multi-Image (train)  & 4823 & 2090 & 5878 \\
        Multi-Image (test)   & 509  & 218  & 636  \\
        \midrule
        Total train & 6096 & 2090 & 18399 \\
        Total test  & 651  & 218  & 2010  \\
        \bottomrule
    \end{tabular}
\end{table}

% -----------------------------------Conclusion-----------------------------------
\section{\textbf{Conclusion}}

% -----------------------------------Reference-----------------------------------
\bibliographystyle{IEEEtran}
\bibliography{ref}

\end{document}
